\section{Resultados esperados}

Se espera obtener, en primer lugar, un marco de simulación reproducible capaz de ejecutar escenarios comparables de transporte cooperativo multi-AGV bajo comunicación local. Ese resultado incluye código organizado, configuraciones de experimento trazables, métricas definidas de forma explícita y un conjunto mínimo de figuras y tablas útiles para la memoria del TFM.

En segundo lugar, se espera producir una comparación técnicamente clara entre la línea base nominal y la variante adaptativa o robusta. El resultado valioso no consiste necesariamente en demostrar superioridad absoluta de la segunda, sino en identificar con precisión en qué escenarios mejora el desempeño, cuándo esa mejora es marginal y qué coste adicional introduce en esfuerzo de control o complejidad de ajuste.

En tercer lugar, se espera cerrar una conclusión académicamente defendible sobre la viabilidad del enfoque propuesto dentro del alcance adoptado. Incluso un resultado negativo o mixto sería útil si permite sostener, con evidencia cuantitativa, que ciertas formas de incertidumbre afectan más que otras o que la ganancia en robustez no compensa determinados costes.

Como productos concretos del TFM, se prevén:
\begin{itemize}
  \item una propuesta de modelado 2D adecuada al problema;
  \item un controlador nominal y una extensión adaptativa o robusta implementados en simulación;
  \item una campaña experimental documentada;
  \item una memoria con discusión crítica de hallazgos, límites y líneas futuras.
\end{itemize}
